\section{Introduction}\label{sec:intro}

A filesystem trusts the bytes a read brings back. Between the medium and
the page cache those bytes cross a controller, a bus, DMA buffers and
caches, any of which can flip a bit: soft errors in commercial silicon
are a measured fact~\cite{baumann2005}, and they matter most where
hardware is commercial and the environment is not. Memon et
al.~\cite{memon2025tr,memon2025arxiv} irradiated commercial systems on
chip running Linux with 20 to 58\,MeV protons and found that, on an
i.MX~8M~Plus running ext4 on eMMC, filesystem faults and storage-driver
errors accounted for 90\,\% of the observed failures, up to journal
aborts and read-only remounts~\cite{memon2025arxiv}. Conventional
filesystems either do not check file data (ext4, ext3) or check it and,
holding a single copy, refuse the read (btrfs); they return the data the
application wrote only from a second copy of
it~\cite{btrfsautorepair,zfschecksums}.

\bfs~\cite{beamfsimpl} is a Linux filesystem built to return it, or to
refuse. Every data block is a \emph{capsule} of sixteen interleaved
RS(255,239) codewords~\cite{reed1960,macwilliams1977}; every read decodes
it; a block beyond correction fails the read with an error instead of
handing back wrong bytes. \bfs descends from FTRFS~\cite{fuchs2015ftrfs}
through an RFC posted in April 2026~\cite{ftrfsrfc} and two earlier
reports~\cite{desbrieres2026ftrfs,desbrieres2026beamfsv2}.
\Cref{sec:ftrfs} explains why it is a new filesystem rather than a
revision of FTRFS.

\paragraph*{What this report adds}
\begin{enumerate}
\item A measured \bfs 0.1.26: of the \xfsTotal{} tests of the xfstests
  \texttt{auto} group, all \xfsXRan{} that ran on x86-64
  passed, and \xfsAPass{} of the \xfsARan{} that ran on aarch64; the
  tests xfstests declined, mostly for features \bfs does not implement,
  are counted and classified, not passed (\cref{sec:xfstests}).
\item A comparison with ext4, ext3, btrfs and squashfs under emulated
  single-bit upsets on the read path, replicated over three runs, and a
  four-node cluster campaign (\cref{sec:multifs,sec:cluster}).
\item An audit of the injector. emufi 0.8.1~\cite{desbrieres2026emufi}
  logs one of its two read hooks one bio early on this kernel. We locate
  the cause in the kernel source, reconstruct the hook of every flip,
  and check the reconstruction against the kernel's own correction log
  (\cref{sec:injector}).
\item A fail-closed event on an indirect block, explained from the code,
  and the residual it exposes (\cref{sec:indirect}).
\item A parallel with the proton campaign of Memon et al., including what
  that work reports about the hardware it irradiated (\cref{sec:memon}).
\item A positioning of \bfs against redundancy, with its costs and
  where it is weaker today (\cref{sec:brings}), and the work the next
  version takes on (\cref{sec:v4}).
\item A reproducibility record in which every number of the text is
  generated from GPG-signed run manifests and frozen by hash
  (\apprepro).
\end{enumerate}

\paragraph*{What changed since v2}
The v2 report~\cite{desbrieres2026beamfsv2} is superseded on five
points. (i)~Data blocks are now interleaved capsules
(\cref{sec:capsule}); v2 described the alternating layout. (ii)~v2 stated
that a corrected block is written back to the medium on read; in
0.1.26 the read path corrects into a private copy and does not write
back, on-disk repair being left to the background scrubber
(\cref{sec:readpath}). (iii)~The module is no longer 2\,198 lines: it is
\moduleLines{} lines of C and headers at \texttt{\commitBeamfs}, beyond the
5\,000-line audit target v2 quoted. (iv)~Injection moved from RadFI~\cite{desbrieres2026radfi} to
emufi. (v)~Validation is the xfstests \texttt{auto} group, test by test,
on two architectures.
